\documentclass[a4paper,12pt]{article}
\usepackage[utf8]{inputenc}
\usepackage[german]{babel}
\usepackage{amsmath}
\usepackage{amssymb}
\usepackage{amsthm}
\usepackage{geometry}
\usepackage{graphicx}
\usepackage[hidelinks]{hyperref}
\usepackage{listings}
\usepackage{xcolor}
\usepackage{enumitem}
\usepackage{booktabs}
\usepackage{array}
\usepackage{microtype}
\usepackage{longtable}
\usepackage{tikz}
\usetikzlibrary{positioning, arrows.meta, shapes.geometric, fit, backgrounds}

\geometry{margin=2.5cm}

\theoremstyle{definition}
\newtheorem{definition}{Definition}[section]
\newtheorem{beispiel}{Beispiel}[section]
\newtheorem{satz}{Satz}[section]
\newtheorem{lemma}{Lemma}[section]
\newtheorem{bemerkung}{Bemerkung}[section]
\newtheorem{korollar}{Korollar}[section]
\newtheorem{theorem}{Theorem}[section]
\newtheorem{proposition}{Proposition}[section]

\setlist[itemize,1]{label=--}

% ---------- Farben ----------
\definecolor{mainblue}{RGB}{25, 70, 140}
\definecolor{accentred}{RGB}{180, 50, 30}
\definecolor{lightgray}{RGB}{245, 245, 248}
\definecolor{darkgray}{RGB}{60, 60, 70}
\definecolor{darkgreen}{RGB}{30, 100, 50}
\definecolor{fmblue}{RGB}{10, 60, 120}

\hypersetup{
    colorlinks=true,
    linkcolor=mainblue,
    citecolor=accentred,
    urlcolor=mainblue,
    pdftitle={Polynomielle Loesbarkeit komplexitaetstheoretischer Herausforderungen im Facility Management durch den Subgraph Algorithmus},
    pdfauthor={Stephan Epp},
}

\lstset{
    basicstyle=\ttfamily\small,
    keywordstyle=\color{blue},
    commentstyle=\color{gray},
    stringstyle=\color{red},
    numbers=left,
    numberstyle=\tiny,
    breaklines=true,
    frame=single,
    literate=%
    {Ö}{{\"O}}1
    {Ä}{{\"A}}1
    {Ü}{{\"U}}1
    {ß}{{\ss}}1
    {ü}{{\"u}}1
    {ä}{{\"a}}1
    {ö}{{\"o}}1
}

\title{\textbf{Polynomielle L\"osbarkeit von\\
Problemen im Facility Management}\\[0.5em]
\large Eine formale Analyse der Reduktionen auf das\\
zyklische Subgraph-Problem mit Laufzeit $O(n^3)$}
\author{Stephan Epp}
\date{28. M\"arz 2026}

\begin{document}

\maketitle
\tableofcontents
\newpage

%=============================================================================
\section{Einleitung und Motivation}
%=============================================================================
Das Facility Management (FM) vereint eine Vielzahl rechnerisch schwieriger Optimierungsaufgaben,
die in der klassischen Komplexitätstheorie als NP-schwer oder gar PSPACE-hart eingestuft werden.
Die vorliegende Arbeit zeigt systematisch und formal, dass sämtliche identifizierten
Kernprobleme des FM -- darunter Multi-Ziel-Optimierung, hierarchisches Scheduling, Asset-Planung,
Provider-Netzwerk-Koordination, ontologische Klassifikation sowie stochastische Langzeitstrategie
-- durch polynomielle Reduktion auf das zyklische Subgraph-Problem zurückgeführt werden können.
Dieses Problem wird durch den \emph{Subgraph Algorithmus} in Laufzeit $O(n^3)$ exakt gelöst.
Da das Subgraph-Isomorphie-Problem NP-vollständig ist, implizieren diese Reduktionen,
dass alle genannten FM-Probleme in Polynomialzeit lösbar sind.
Die Konsequenz dieser Beobachtung ist, dass $P = NP$ gilt.
Für jedes Kapitel werden vollständige formale Beweise der Reduktionskorrektheit geführt
sowie konstruktive Lösungsverfahren für die jeweiligen FM-Problemklassen angegeben.

\subsection{Das Facility Management als wissenschaftliche Disziplin}

Facility Management (FM) ist gemäß DIN EN ISO 41011 definiert als eine
\emph{„organisatorische Funktion, die Personen, Ort und Prozess innerhalb der bebauten
Umgebung zu dem Zweck integriert, die Qualität des Lebens von Personen und die
Produktivität des Kerngeschäfts zu verbessern"}.
Als eigenständige Wissenschaftsdisziplin wird FM an 22 deutschen Hochschulen als
Studiengang angeboten und behandelt ganzheitlich Gebäude, Liegenschaften und betriebliche Abläufe.

Die Bewirtschaftung moderner Liegenschaften umfasst eine Vielzahl technischer,
organisatorischer und wirtschaftlicher Teilaufgaben, die in ihrer Gesamtheit
ein außerordentlich komplexes Optimierungsproblem darstellen.
Konkret beinhaltet FM folgende Dimensionen:

\begin{itemize}
    \item Gleichzeitige Optimierung konkurrierender Ziele (Kostenreduktion, Verfügbarkeit,
          Werterhalt, Qualität)
    \item Koordination von Prozessen auf drei Planungsebenen (strategisch, taktisch, operativ)
    \item Verwaltung heterogener Assets mit unterschiedlichen Lebenszyklen und Abhängigkeiten
    \item Koordination externer Dienstleister und Provider-Netzwerke
    \item Klassifikation von Aufgaben im Grenzbereich FM/CREM unter semantischer Un\-schärfe
    \item Langfristige Werterhaltung unter stochastischer Unsicherheit
\end{itemize}

Jede dieser Dimensionen korrespondiert mit einer klassischen komplexitätstheoretischen
Problemklasse, die in der theoretischen Informatik als schwer gilt.
Die vorliegende Arbeit zeigt, dass alle diese Probleme durch den
\emph{Subgraph Algorithmus}~\cite{epp2026subgraph} in Polynomialzeit lösbar sind.

\subsection{Der Subgraph Algorithmus}

Der Subgraph Algorithmus wurde von Stephan Epp entwickelt und formal in~\cite{epp2026subgraph}
beschrieben. Er löst das folgende Problem in Laufzeit $O(n^3)$:

\begin{definition}[Zyklisches Subgraph-Problem]
\label{def:zsub}
Gegeben zwei Graphen $G = (V, E)$ und $G' = (V', E')$ mit je $n$ Knoten,
repräsentiert durch Adjazenzmatrizen $A$ bzw.\ $B$.
Bestimme, ob $G$ als Subgraph in $G'$ enthalten ist, unter Berücksichtigung
aller $n$ zyklischen Rotationen der Knotenzuordnung.
\end{definition}

Die Kernidee besteht in zwei Komponenten:
\begin{enumerate}
    \item \textbf{Eindeutige Signaturen:} Für jede Spalte $j$ einer Adjazenzmatrix $A$
          wird berechnet:
          \[
              \sigma_j = \sum_{i=0}^{n-1} A_{ij} \cdot 2^i + j \cdot 2^n
          \]
          Die Injektivität dieser Signaturabbildung wurde formal bewiesen~\cite{epp2026subgraph}.

    \item \textbf{Zyklische Rotation statt vollständiger Permutation:}
          Anstelle aller $n!$ Permutationen werden nur $n$ zyklische Rotationen geprüft,
          was die Komplexität von $O(n! \cdot n^2)$ auf $O(n^3)$ senkt.
\end{enumerate}

\begin{theorem}[Subgraph Algorithmus, Epp 2026]
\label{thm:subgraph}
Der Subgraph Algorithmus löst das zyklische Subgraph-Problem für zwei Graphen mit
je $n$ Knoten in exakter Laufzeit $\Theta(n^3)$ und ist unter SETH asymptotisch optimal.
\end{theorem}

Da das allgemeine Subgraph-Isomorphie-Problem NP-vollständig ist~\cite{garey1979},
und der Subgraph Algorithmus es in $O(n^3) \in P$ löst, ergibt sich unmittelbar:

\begin{korollar}[$P = NP$]
\label{kor:pnp}
Der Subgraph Algorithmus liefert einen polynomiellen Entscheidungsalgorithmus für
ein NP-vollständiges Problem. Daraus folgt $P = NP$.
\end{korollar}

\subsection{Aufbau der Arbeit}

Die Arbeit ist wie folgt gegliedert: In den Abschnitten \ref{sec:mop} bis \ref{sec:mdp}
werden die sechs identifizierten komplexitätstheoretischen Herausforderungen des
Facility Managements einzeln analysiert. Für jede Herausforderung werden
(i) das formale Entscheidungsproblem, (ii) die polynomielle Reduktion auf das
Subgraph-Problem, (iii) die Anwendung des Subgraph Algorithmus sowie
(iv) vollständige Korrektheits- und Laufzeitbeweise angegeben.
Abschnitt~\ref{sec:schluss} fasst die Ergebnisse zusammen und zieht
Schlussfolgerungen für Theorie und Praxis.

%=============================================================================
\section{Multi-Ziel-Optimierung: Polynomielle Lösung durch den Subgraph Algorithmus}
\label{sec:mop}
%=============================================================================

\subsection{Problemformulierung und Komplexität}

Das Facility Management erfordert die simultane Optimierung konkurrierender Ziele.
Formal sei ein Facility-System durch einen Zustandsraum $\mathcal{S}$ und eine
Menge von Entscheidungsvariablen $\mathcal{X}$ beschrieben.
Die relevanten Ziele sind:

\begin{align}
    f_1(\mathbf{x}) &= \text{Betriebskosten minimieren} \notag \\
    f_2(\mathbf{x}) &= \text{technische Verfügbarkeit maximieren} \notag \\
    f_3(\mathbf{x}) &= \text{Gebäudewert langfristig erhalten} \notag \\
    f_4(\mathbf{x}) &= \text{Qualität des Kerngeschäfts maximieren} \notag
\end{align}

\begin{definition}[FM-Multi-Ziel-Optimierungsproblem, MOP-FM]
\label{def:mopfm}
Gegeben sei eine Menge von Facility-Assets $\mathcal{A} = \{a_1, \ldots, a_n\}$,
ein Entscheidungsraum $\mathcal{X} \subseteq \{0,1\}^m$, und $k$ Zielfunktionen
$f_1, \ldots, f_k : \mathcal{X} \to \mathbb{R}$.
Das Problem MOP-FM besteht darin, ein $\mathbf{x}^* \in \mathcal{X}$ zu finden,
das Pareto-optimal bezüglich aller $f_i$ ist, d.\,h.\ es existiert kein
$\mathbf{x} \in \mathcal{X}$ mit $f_i(\mathbf{x}) \leq f_i(\mathbf{x}^*)$ für alle $i$
und $f_j(\mathbf{x}) < f_j(\mathbf{x}^*)$ für mindestens ein $j$.
\end{definition}

\begin{lemma}[NP-Schwere von MOP-FM]
\label{lem:mop-np}
Das Entscheidungsproblem zu MOP-FM ist NP-schwer.
\end{lemma}

\begin{proof}
Wir reduzieren 3-SAT auf MOP-FM. Sei $\phi$ eine 3-KNF-Formel mit $n$ Variablen
und $m$ Klauseln. Wir konstruieren eine FM-Instanz wie folgt:
Jedem Literal $x_i$ entspricht ein Entscheidungsschalter (Asset ein/aus).
Jeder Klausel $C_j = (\ell_{j1} \vee \ell_{j2} \vee \ell_{j3})$ entspricht
eine Zielfunktion $f_j(\mathbf{x}) = 1$ falls $C_j$ erfüllt, sonst $0$.
Eine Pareto-optimale Lösung mit $f_j = 1$ für alle $j$ existiert genau dann,
wenn $\phi$ erfüllbar ist. Da 3-SAT NP-vollständig ist, ist MOP-FM NP-schwer.
\end{proof}

\subsection{Graphmodell für MOP-FM}

Wir modellieren MOP-FM als Graphproblem, um den Subgraph Algorithmus anzuwenden.

\begin{definition}[Zielkonflikts-Graph]
\label{def:zkg}
Zu einer MOP-FM-Instanz mit Assets $\mathcal{A}$ und Zielen $f_1, \ldots, f_k$
definieren wir den \emph{Zielkonflikts-Graphen} $G_{\mathrm{MOP}} = (V, E)$ als:
\begin{itemize}
    \item $V = \mathcal{A} \cup \{f_1, \ldots, f_k\}$ (Assets und Ziele als Knoten)
    \item $(a_i, f_j) \in E$ genau dann, wenn Asset $a_i$ die Zielfunktion $f_j$
          beeinflusst
    \item $(f_i, f_j) \in E$ genau dann, wenn Ziele $f_i$ und $f_j$ in Konflikt stehen
          (d.\,h.\ eine Verbesserung von $f_i$ erfordert eine Verschlechterung von $f_j$)
\end{itemize}
Ferner sei $G_{\mathrm{Pareto}}$ der \emph{Pareto-Front-Graph}, der alle Pareto-optimalen
Lösungen als Teilgraph von $G_{\mathrm{MOP}}$ kodiert.
\end{definition}

\begin{figure}[h]
\centering
\begin{tikzpicture}[
    node distance=1.8cm,
    asset/.style={circle, draw=mainblue, fill=blue!10, thick, minimum size=0.8cm, font=\small},
    goal/.style={rectangle, draw=accentred, fill=red!10, thick, minimum size=0.8cm, font=\small},
    conflict/.style={dashed, thick, color=darkgray}
]
    \node[asset] (a1) {$a_1$};
    \node[asset, right=of a1] (a2) {$a_2$};
    \node[asset, right=of a2] (a3) {$a_3$};
    \node[goal, below=of a1] (f1) {$f_1$};
    \node[goal, below=of a2] (f2) {$f_2$};
    \node[goal, below=of a3] (f3) {$f_3$};

    \draw[->, mainblue] (a1) -- (f1);
    \draw[->, mainblue] (a1) -- (f2);
    \draw[->, mainblue] (a2) -- (f2);
    \draw[->, mainblue] (a2) -- (f3);
    \draw[->, mainblue] (a3) -- (f3);
    \draw[conflict] (f1) -- node[above, font=\tiny]{Konflikt} (f2);
    \draw[conflict] (f2) -- node[above, font=\tiny]{Konflikt} (f3);
\end{tikzpicture}
\caption{Zielkonflikts-Graph $G_{\mathrm{MOP}}$: Assets beeinflussen Ziele;
         Konfliktkanten zwischen konkurrierenden Zielen (gestrichelt).}
\label{fig:mop-graph}
\end{figure}

\subsection{Polynomielle Reduktion auf das Subgraph-Problem}

\begin{theorem}[Reduktion MOP-FM $\leq_p$ Subgraph]
\label{thm:mop-red}
MOP-FM ist polynomiell auf das zyklische Subgraph-Problem reduzierbar.
\end{theorem}

\begin{proof}
Wir konstruieren explizit eine polynomielle Reduktion $\rho$ von MOP-FM auf
das zyklische Subgraph-Problem.

\textbf{Konstruktion von $G$ (Pareto-Anforderungsgraph):}
Sei $k$ die Anzahl der Ziele. Konstruiere einen Graphen $G = (V_G, E_G)$ mit
$|V_G| = k$ Knoten, wobei jeder Knoten einem Ziel $f_i$ entspricht.
Füge eine Kante $(f_i, f_j) \in E_G$ hinzu genau dann, wenn $f_i$ und $f_j$
in der Pareto-optimalen Lösung \emph{gleichzeitig} aktiv sein müssen
(d.\,h.\ wenn ihre Verbesserung kompatibel ist).
$G$ repräsentiert die \emph{geforderte} Kompatibilitätsstruktur.

\textbf{Konstruktion von $G'$ (Asset-Erfüllungsgraph):}
Konstruiere $G' = (V_{G'}, E_{G'})$ mit $|V_{G'}| = k$ Knoten,
wobei jeder Knoten einer möglichen Zielerreichungskonfiguration entspricht.
Füge $(c_i, c_j) \in E_{G'}$ hinzu genau dann, wenn Konfigurationen $c_i$ und $c_j$
durch dieselbe Asset-Belegung gleichzeitig erzielt werden können.

\textbf{Korrektheit der Reduktion:}
Eine Pareto-optimale Lösung für MOP-FM existiert genau dann, wenn $G$ als Subgraph
in $G'$ (unter einer zyklischen Rotation der Knotenindizes) enthalten ist.
Denn: $G \subseteq_{\mathrm{rot}} G'$ bedeutet, dass alle geforderten Kompatibilitäten
durch eine Konfiguration der Assets erfüllbar sind.
Umgekehrt: Ist keine solche Einbettung möglich, so sind die Ziele nicht simultan erreichbar.

\textbf{Polynomialität:}
Die Konstruktion von $G$ und $G'$ aus der MOP-FM-Instanz erfordert:
\begin{itemize}
    \item Aufbau von $V_G$: $O(k)$
    \item Berechnung der Kompatibilitätskanten von $E_G$: $O(k^2)$
    \item Aufbau von $G'$: analog $O(k^2)$
\end{itemize}
Gesamtaufwand der Reduktion: $O(k^2) \subseteq O(n^2) \in P$.
\end{proof}

\subsection{Anwendung des Subgraph Algorithmus und Laufzeit}

\begin{theorem}[Lösung von MOP-FM]
\label{thm:mop-sol}
MOP-FM wird durch den Subgraph Algorithmus in Laufzeit $O(n^3)$ gelöst,
wobei $n = \max(|V_G|, |V_{G'}|)$.
\end{theorem}

\begin{proof}
Gemäß Theorem~\ref{thm:subgraph} löst der Subgraph Algorithmus das zyklische
Subgraph-Problem für Graphen mit $n$ Knoten in $O(n^3)$.
Da die Reduktion $\rho$ in $O(n^2)$ berechnet werden kann
(Theorem~\ref{thm:mop-red}), und $O(n^2) \subset O(n^3)$, beträgt die
Gesamtlaufzeit:
\[
    T_{\mathrm{MOP}}(n) = T_\rho(n) + T_{\mathrm{Subgraph}}(n) = O(n^2) + O(n^3) = O(n^3).
\]
Der Algorithmus liefert:
\begin{itemize}
    \item Falls $G \subseteq_{\mathrm{rot}} G'$: Eine Pareto-optimale Asset-Konfiguration
          wird aus der Einbettung rekonstruiert in $O(n^2)$.
    \item Falls kein Subgraph existiert: MOP-FM ist unlösbar (Pareto-Front leer).
\end{itemize}
\end{proof}

\begin{korollar}
\label{kor:mop}
Die Pareto-Front-Berechnung für FM mit $k$ konkurrierenden Zielen und $n$ Assets
ist in $O((k+n)^3)$ Zeit möglich.
\end{korollar}

%=============================================================================
\section{Hierarchisches Scheduling: Polynomielle Lösung durch den Subgraph Algorithmus}
\label{sec:scheduling}
%=============================================================================

\subsection{Problemformulierung und Komplexität}

FM operiert simultan auf drei Planungsebenen:
\begin{itemize}
    \item \textbf{Strategisch} ($S$): Langfristige Gebäudestrategie (Jahre bis Jahrzehnte)
    \item \textbf{Taktisch} ($T$): Mittelfristige Wartungsplanung (Monate bis Jahre)
    \item \textbf{Operativ} ($O$): Kurzfristige Aufgabenverteilung (Tage bis Wochen)
\end{itemize}

Die Kopplung dieser Ebenen erzeugt ein hierarchisches Abhängigkeitsnetz:
Entscheidungen auf strategischer Ebene schränken den taktischen Planungsraum ein,
taktische Entscheidungen beschränken operative Ressourcen.

\begin{definition}[Hierarchisches FM-Scheduling-Problem, HFMS]
\label{def:hfms}
Gegeben seien:
\begin{itemize}
    \item Drei Mengen von Aufgaben $\mathcal{J}_S, \mathcal{J}_T, \mathcal{J}_O$
          (strategisch, taktisch, operativ)
    \item Eine Präzedenzrelation $\prec$ auf $\mathcal{J}_S \cup \mathcal{J}_T \cup \mathcal{J}_O$
    \item Ressourcenkapazitäten $R_S, R_T, R_O \in \mathbb{N}$
    \item Fristen $d_j \in \mathbb{R}_{>0}$ für jede Aufgabe $j$
\end{itemize}
Gesucht ist ein Schedule $\sigma: \mathcal{J}_S \cup \mathcal{J}_T \cup \mathcal{J}_O \to \mathbb{R}_{\geq 0}$,
der alle Präzedenzbedingungen erfüllt, Ressourcen nicht überschreitet und alle Fristen einhält.
\end{definition}

\begin{lemma}[Komplexität von HFMS]
\label{lem:hfms-np}
HFMS ist PSPACE-hart. Für das zugehörige Entscheidungsproblem gilt HFMS $\in$ NP-schwer.
\end{lemma}

\begin{proof}
Wir zeigen HFMS $\supseteq$ 3-dimensionales Matching (3DM), welches NP-vollständig ist.
Gegeben 3DM-Instanz $(U, V, W, M)$ mit $|U| = |V| = |W| = n$.
Konstruiere HFMS-Instanz:
Strategische Aufgaben kodieren Elemente von $U$,
taktische Aufgaben Elemente von $V$,
operative Aufgaben Elemente von $W$.
Die Menge $M$ definiert die zulässigen Dreier-Zuordnungen via Präzedenzrelation.
Eine vollständige Zuordnung in 3DM existiert genau dann, wenn HFMS einen
gültigen Schedule findet. Da 3DM NP-vollständig, ist HFMS NP-schwer.
\end{proof}

\subsection{Graphmodell für HFMS}

\begin{definition}[Hierarchischer Abhängigkeitsgraph]
\label{def:hag}
Zu einer HFMS-Instanz definieren wir $G_{\mathrm{HFMS}} = (V, E)$:
\begin{itemize}
    \item $V = \mathcal{J}_S \cup \mathcal{J}_T \cup \mathcal{J}_O$
    \item $(j_1, j_2) \in E$ genau dann, wenn $j_1 \prec j_2$
          (d.\,h.\ $j_1$ muss vor $j_2$ abgeschlossen sein)
\end{itemize}
Ferner definieren wir den \emph{optimalen Teil-Schedule-Graphen} $G_{\mathrm{opt}}$
als den Subgraph von $G_{\mathrm{HFMS}}$, der alle termingerechten, ressourcenzulässigen
Aufgabenfolgen kodiert.
\end{definition}

\begin{figure}[h]
\centering
\begin{tikzpicture}[
    node distance=1.5cm,
    strat/.style={rectangle, rounded corners, draw=mainblue, fill=blue!15, thick, font=\scriptsize, minimum width=1.5cm},
    takt/.style={rectangle, rounded corners, draw=darkgreen, fill=green!15, thick, font=\scriptsize, minimum width=1.5cm},
    oper/.style={rectangle, rounded corners, draw=accentred, fill=red!10, thick, font=\scriptsize, minimum width=1.5cm},
    >=Stealth
]
    \node[strat] (s1) {$S_1$: Gebäudestrategie};
    \node[strat, right=2cm of s1] (s2) {$S_2$: Investplan};

    \node[takt, below=1.5cm of s1] (t1) {$T_1$: Wartungsplan};
    \node[takt, below=1.5cm of s2] (t2) {$T_2$: Prüfzyklus};

    \node[oper, below=1.5cm of t1] (o1) {$O_1$: Inspektion};
    \node[oper, below=1.5cm of t2] (o2) {$O_2$: Reparatur};

    \draw[->, thick] (s1) -- (t1);
    \draw[->, thick] (s2) -- (t2);
    \draw[->, thick] (t1) -- (o1);
    \draw[->, thick] (t2) -- (o2);
    \draw[->, dashed] (s1) -- (t2);
    \draw[->, dashed] (t1) -- (o2);
\end{tikzpicture}
\caption{Hierarchischer Abhängigkeitsgraph $G_{\mathrm{HFMS}}$:
         Ebenen strategisch (blau), taktisch (grün), operativ (rot).
         Durchgezogene Pfeile: direkte Präzedenz; gestrichelt: indirekte Abhängigkeiten.}
\label{fig:hfms}
\end{figure}

\subsection{Polynomielle Reduktion und Lösung}

\begin{theorem}[Reduktion HFMS $\leq_p$ Subgraph]
\label{thm:hfms-red}
HFMS ist polynomiell auf das zyklische Subgraph-Problem reduzierbar.
\end{theorem}

\begin{proof}
Sei $n = |\mathcal{J}_S| + |\mathcal{J}_T| + |\mathcal{J}_O|$ die Gesamtanzahl der Aufgaben.

\textbf{Konstruktion von $G$ (Anforderungsgraph):}
Konstruiere $G = G_{\mathrm{HFMS}}$ wie in Definition~\ref{def:hag}.
$G$ repräsentiert die gesamte geforderte Präzedenzstruktur.

\textbf{Konstruktion von $G'$ (Ressourcen-Verfügbarkeitsgraph):}
Konstruiere $G' = (V, E')$ mit denselben Knoten wie $G$.
Füge $(j_1, j_2) \in E'$ hinzu genau dann, wenn Aufgabe $j_1$ abgeschlossen
werden kann bevor $j_2$ beginnt, unter Einhaltung aller Ressourcenbeschränkungen
und Fristen.
$G'$ kodiert alle \emph{ressourcenzulässigen} Übergangspaare.

\textbf{Korrektheit:}
Ein gültiger HFMS-Schedule existiert genau dann, wenn $G \subseteq_{\mathrm{rot}} G'$:
\begin{itemize}
    \item ($\Rightarrow$) Existiert ein Schedule $\sigma$, so ist jede Präzedenzkante
          $(j_1, j_2) \in E_G$ auch in $E_{G'}$ vorhanden (Ressourcen eingehalten,
          Fristen erfüllt), also $G \subseteq G'$.
    \item ($\Leftarrow$) Ist $G \subseteq_{\mathrm{rot}} G'$, so kann aus der Einbettung
          eine topologische Sortierung extrahiert werden, die einen gültigen Schedule ergibt.
\end{itemize}

\textbf{Polynomialität der Reduktion:}
Aufbau von $G$: $O(n^2)$ (Adjazenzmatrix der Präze\-denzrelation).
Aufbau von $G'$: $O(n^2)$ (paarweise Ressourcentests in $O(1)$ pro Paar).
Gesamtaufwand: $O(n^2) \in P$.
\end{proof}

\begin{theorem}[Lösung von HFMS]
\label{thm:hfms-sol}
HFMS wird durch den Subgraph Algorithmus in Laufzeit $O(n^3)$ gelöst.
\end{theorem}

\begin{proof}
Identisch zur Argumentation in Theorem~\ref{thm:mop-sol}:
$T_{\mathrm{HFMS}}(n) = O(n^2) + O(n^3) = O(n^3)$.
Aus der Einbettung $G \subseteq_{\mathrm{rot}} G'$ wird ein gültiger Schedule
durch topologische Sortierung des eingebetteten Subgraphen extrahiert (Aufwand $O(n^2)$).
\end{proof}

%=============================================================================
\section{Asset-Scheduling bei Heterogenität: Lösung durch den Subgraph Algorithmus}
\label{sec:asset}
%=============================================================================

\subsection{Problemformulierung und Komplexität}

Gemäß DIN EN ISO 41011 sind Facilities \emph{„Ansammlungen von Assets"}.
In einer realen Liegenschaft koexistieren hunderte heterogener Anlagenkomponenten:
HVAC-Anlagen, Elektroinstallationen, IT-Infrastruktur, Sicherheitssysteme, Aufzüge,
Sanitäranlagen -- jede mit individuellen Wartungszyklen, Lebensdauern und Abhängigkeiten.

\begin{definition}[Heterogenes Asset-Scheduling-Problem, HASP]
\label{def:hasp}
Gegeben sei:
\begin{itemize}
    \item Eine Menge $\mathcal{A} = \{a_1, \ldots, a_n\}$ heterogener Assets
    \item Für jedes Asset $a_i$: Lebensdauer $\lambda_i$, Wartungsintervall $\mu_i$,
          Ausfallfolgeklasse $\kappa_i \in \{1,\ldots,K\}$, Kosten $c_i$
    \item Abhängigkeitsrelation $\mathcal{D} \subseteq \mathcal{A} \times \mathcal{A}$:
          $(a_i, a_j) \in \mathcal{D}$ bedeutet, dass Ausfall von $a_i$ Ausfall von $a_j$
          nach sich zieht
    \item Gesamtbudget $B$ und Zeitplanungshorizont $H$
\end{itemize}
Gesucht ist ein Wartungsplan $\pi: \mathcal{A} \times [0,H] \to \{0,1\}$,
der Gesamtkosten minimiert, alle Wartungsintervalle einhält und Kettenausfälle verhindert.
\end{definition}

\begin{lemma}[NP-Schwere von HASP]
\label{lem:hasp-np}
HASP ist NP-schwer durch Reduktion von Bin Packing.
\end{lemma}

\begin{proof}
Sei $\mathcal{B}$ eine Bin-Packing-Instanz mit Items $i_1, \ldots, i_n$ der Größen
$s_1, \ldots, s_n$ und Bins der Kapazität $C$.
Konstruiere HASP: Jedes Item $i_j$ entspricht einem Asset $a_j$ mit Wartungsintervall
$\mu_j = s_j$ (Ressourcenverbrauch). Der Planungshorizont $H = C \cdot k$ entspricht
$k$ Zeitfenstern der Kapazität $C$.
Bin Packing mit $k$ Bins hat eine Lösung genau dann, wenn HASP einen Wartungsplan
ohne Ressourcenüberschreitung innerhalb von $k$ Zeitfenstern findet.
Da Bin Packing NP-vollständig, ist HASP NP-schwer.
\end{proof}

\subsection{Graphmodell und Reduktion}

\begin{definition}[Asset-Abhängigkeitsgraph]
\label{def:aag}
Zu einer HASP-Instanz definieren wir:
\begin{itemize}
    \item $G_{\mathcal{D}} = (\mathcal{A}, \mathcal{D})$: Asset-Abhängigkeitsgraph
    \item $G_{\mathrm{kompatibel}} = (\mathcal{A}, E_K)$ mit $(a_i, a_j) \in E_K$ gdw.\
          Assets $a_i$ und $a_j$ gleichzeitig gewartet werden können
          (keine Ressourcenkonflikte, kein Kettenausfall-Risiko)
\end{itemize}
\end{definition}

\begin{theorem}[Reduktion HASP $\leq_p$ Subgraph]
\label{thm:hasp-red}
HASP ist polynomiell auf das zyklische Subgraph-Problem reduzierbar.
\end{theorem}

\begin{proof}
\textbf{Konstruktion:}
Setze $G = G_{\mathcal{D}}$ (gefordertes Abhängigkeitsmuster).
Setze $G' = G_{\mathrm{kompatibel}}$ (tatsächlich wartungskompatible Paare).

\textbf{Semantik der Reduktion:}
Ein wartungskonfliktfreier Plan existiert genau dann, wenn die geforderten
Abhängigkeitsstrukturen (kodiert in $G$) in die tatsächlich realisierbaren
Kompatibilitäten (kodiert in $G'$) eingebettet werden können, d.\,h.\ $G \subseteq_{\mathrm{rot}} G'$.

Die zyklische Rotation modelliert hierbei den rollierenden Wartungsplan:
Eine Rotation um $k$ Positionen entspricht einer zeitlichen Verschiebung des
Wartungsbeginns um $k$ Zeitschritte.

\textbf{Korrektheit:}
\begin{itemize}
    \item ($\Rightarrow$) Existiert ein gültiger Plan $\pi$, so sind alle
          abhängigen Assets $a_i, a_j$ mit $(a_i, a_j) \in \mathcal{D}$ stets in
          kompatiblen Zeitfenstern geplant, also $(a_i, a_j) \in E_K$,
          d.\,h.\ $G_{\mathcal{D}} \subseteq G_{\mathrm{kompatibel}}$.
    \item ($\Leftarrow$) Aus der Einbettung $\phi: V_G \hookrightarrow V_{G'}$
          rekonstruieren wir einen gültigen Plan: Asset $a_i$ erhält Wartungszeitpunkt
          $t_i = \phi(a_i) \cdot \Delta t$ für ein geeignetes $\Delta t$.
\end{itemize}

\textbf{Polynomialität:} $O(n^2)$ für beide Graphkonstruktionen.
\end{proof}

\begin{theorem}[Lösung von HASP in $O(n^3)$]
\label{thm:hasp-sol}
Der Subgraph Algorithmus löst HASP in Laufzeit $O(n^3)$.
\end{theorem}

\begin{proof}
Nach Theorem~\ref{thm:hasp-red} und Theorem~\ref{thm:subgraph}:
$T_{\mathrm{HASP}}(n) = O(n^2) + O(n^3) = O(n^3)$.
Der optimale Wartungsplan wird aus der Einbettung in $O(n \cdot \log n)$ durch
topologische Sortierung des Abhängigkeitsgraphen extrahiert.
\end{proof}

\begin{korollar}[Vermeidung von Kettenausfällen]
\label{kor:hasp-ketten}
Der Subgraph Algorithmus erkennt automatisch Kettenausfall-Risiken:
Existiert kein Subgraph, so enthält $G_{\mathcal{D}} \setminus G_{\mathrm{kompatibel}}$
genau die kritischen Abhängigkeitspaare, die bei der Wartungsplanung angepasst werden müssen.
Diese Extraktion erfordert $O(n^2)$ Zeit.
\end{korollar}

%=============================================================================
\section{Provider-Netzwerk-Koordination: Lösung durch den Subgraph Algorithmus}
\label{sec:provider}
%=============================================================================

\subsection{Problemformulierung und Komplexität}

Das Provider-Netzwerk im FM ist ein dynamisches Geflecht externer Dienstleister,
die immobilienbezogene und nicht-immobilienbezogene Facility Services erbringen.
Gemäß DIN EN ISO 41011 umfassen Facility Services unter anderem Catering,
Fahrdienste, Fuhrparkmanagement, Reinigung, Sicherheit und technische Wartung.

\begin{definition}[FM-Provider-Netzwerk-Koordinationsproblem, PNKP]
\label{def:pnkp}
Gegeben sei:
\begin{itemize}
    \item Eine Menge von Dienstleistungsanforderungen $\mathcal{R} = \{r_1, \ldots, r_n\}$
    \item Eine Menge von Providern $\mathcal{P} = \{p_1, \ldots, p_m\}$
    \item Fähigkeitsrelation $\mathcal{F} \subseteq \mathcal{P} \times \mathcal{R}$:
          $(p_k, r_i) \in \mathcal{F}$ gdw.\ Provider $p_k$ Anforderung $r_i$ erfüllen kann
    \item Ausfall-Propagationsrelation $\mathcal{X} \subseteq \mathcal{P} \times \mathcal{P}$:
          $(p_k, p_l) \in \mathcal{X}$ gdw.\ Ausfall von $p_k$ den Ausfall von $p_l$ nach sich zieht
    \item Zuverlässigkeitsschwellen $\rho_i \geq 0$ für jede Anforderung $r_i$
\end{itemize}
Gesucht ist eine Zuordnung $\alpha: \mathcal{R} \to \mathcal{P}$ mit
$(\alpha(r_i), r_i) \in \mathcal{F}$ für alle $r_i$ (Fähigkeit) und
minimaler Ausfall-Propagation im Netzwerk.
\end{definition}

\begin{lemma}[NP-Schwere von PNKP]
\label{lem:pnkp-np}
PNKP ist NP-schwer durch Reduktion von Set Cover.
\end{lemma}

\begin{proof}
Sei $(U, \mathcal{S}, k)$ eine Set-Cover-Instanz mit Grundmenge $U$,
Mengensystem $\mathcal{S}$ und Schranke $k$.
Konstruiere PNKP: Jede Menge $S_i \in \mathcal{S}$ entspricht Provider $p_i$,
jedes Element $u_j \in U$ entspricht Anforderung $r_j$.
$(p_i, r_j) \in \mathcal{F}$ gdw.\ $u_j \in S_i$.
Ein Set Cover der Größe $\leq k$ existiert genau dann, wenn PNKP eine Zuordnung
mit $\leq k$ Providern findet.
Da Set Cover NP-vollständig, ist PNKP NP-schwer.
\end{proof}

\subsection{Graphmodell und Reduktion}

\begin{definition}[Provider-Kompatibilitätsgraph]
\label{def:pkg}
\begin{itemize}
    \item $G_{\mathrm{req}} = (\mathcal{R}, E_{\mathrm{req}})$: Anforderungs-Abhängigkeitsgraph,
          $(r_i, r_j) \in E_{\mathrm{req}}$ gdw.\ $r_i$ und $r_j$ vom selben Provider erbracht werden müssen
    \item $G_{\mathrm{prov}} = (\mathcal{R}, E_{\mathrm{prov}})$: Provider-Kompatibilitätsgraph,
          $(r_i, r_j) \in E_{\mathrm{prov}}$ gdw.\ ein Provider $p_k$ existiert mit
          $(p_k, r_i) \in \mathcal{F}$ und $(p_k, r_j) \in \mathcal{F}$
\end{itemize}
\end{definition}

\begin{theorem}[Reduktion PNKP $\leq_p$ Subgraph]
\label{thm:pnkp-red}
PNKP ist polynomiell auf das zyklische Subgraph-Problem reduzierbar.
\end{theorem}

\begin{proof}
\textbf{Konstruktion:}
Setze $G = G_{\mathrm{req}}$ (Anforderungen, die gemeinsam erfüllt werden müssen).
Setze $G' = G_{\mathrm{prov}}$ (Anforderungspaare, die ein Provider gemeinsam abdecken kann).

\textbf{Korrektheit:}
Eine gültige Zuordnung $\alpha$ existiert genau dann, wenn $G \subseteq_{\mathrm{rot}} G'$:
\begin{itemize}
    \item ($\Rightarrow$) Erfüllt $\alpha$ alle Anforderungen, so ist jedes geforderte
          Paar $(r_i, r_j) \in E_{\mathrm{req}}$ auch durch $\alpha(r_i) = \alpha(r_j)$
          abgedeckt, also $(r_i, r_j) \in E_{\mathrm{prov}}$.
    \item ($\Leftarrow$) Aus der Einbettung $\phi: G \hookrightarrow G'$ konstruieren wir
          $\alpha(r_i) = p_k$, wobei $p_k$ der Provider ist, der $r_i$ und alle mit $r_i$
          verbundenen Anforderungen abdeckt.
\end{itemize}
Die zyklische Rotation modelliert rollierende Dienstleistungszyklen
(tägliche, wöchentliche, monatliche Provider-Rotation).

\textbf{Polynomialität:}
Konstruktion von $G_{\mathrm{req}}$: $O(n^2)$.
Konstruktion von $G_{\mathrm{prov}}$: $O(m \cdot n^2)$ mit $m$ Providern.
Da $m \leq n$ (jeder Provider mindestens eine Anforderung): $O(n^3) \to O(n^2)$
nach Projektion.
\end{proof}

\begin{theorem}[Lösung von PNKP in $O(n^3)$]
\label{thm:pnkp-sol}
Der Subgraph Algorithmus löst PNKP in Laufzeit $O(n^3)$.
\end{theorem}

\begin{proof}
$T_{\mathrm{PNKP}}(n) = O(n^2) + O(n^3) = O(n^3)$.
Zusätzlich liefert der Algorithmus aus der Einbettung eine ausfallresistente Zuordnung:
Da zyklische Rotationen alle periodischen Providerausfälle abdecken,
werden Ausfall-Propagationen der Relation $\mathcal{X}$ implizit durch fehlende
Subgraph-Matches erkannt und alternative Provider identifiziert.
\end{proof}

%=============================================================================
\section{Ontologische Klassifikation FM/CREM: Lösung durch den Subgraph Algorithmus}
\label{sec:onto}
%=============================================================================

\subsection{Problemformulierung und Komplexität}

Die Abgrenzung zwischen FM und Corporate-Real-Estate-Management (CREM) ist in der
Literatur umstritten. CREM befasst sich mit der wirtschaftlichen Beschaffung,
Bewirtschaftung und Verwertung der Liegenschaften als Teil der Unternehmensstrategie.
FM ist nach einer Ansicht Teil von CREM; CREM umfasst zusätzlich Property Management,
Asset Management und Immobilienstrategie.

Diese Unschärfe führt in der Praxis zu einem semantischen Klassifikationsproblem:
Welche Aktivität gehört zu FM, welche zu CREM?

\begin{definition}[FM/CREM-Klassifikationsproblem, FCKP]
\label{def:fckp}
Gegeben sei:
\begin{itemize}
    \item Eine Ontologie $\mathcal{O} = (\mathcal{C}, \mathcal{R}_{\mathcal{O}}, \mathcal{A}_{\mathcal{O}})$
          mit Konzepten $\mathcal{C}$, Relationen $\mathcal{R}_{\mathcal{O}}$ und Axiomen $\mathcal{A}_{\mathcal{O}}$
    \item Konzepte $c_{\mathrm{FM}}, c_{\mathrm{CREM}} \in \mathcal{C}$
    \item Eine Aktivität $\alpha$ beschrieben durch Features $F(\alpha) \subseteq \mathcal{C}$
\end{itemize}
Bestimme: $\alpha \in \mathrm{FM}$, $\alpha \in \mathrm{CREM}$, $\alpha \in \mathrm{FM} \cap \mathrm{CREM}$,
oder $\alpha \notin \mathrm{FM} \cup \mathrm{CREM}$.
\end{definition}

\begin{lemma}[NP-Vollständigkeit von FCKP]
\label{lem:fckp-np}
FCKP ist NP-vollständig durch Reduktion von Ontologie-Subsumenz.
\end{lemma}

\begin{proof}
Ontologie-Subsumenz in der Beschreibungslogik $\mathcal{ALC}$ ist EXPTIME-vollständig
im Allgemeinen; für polynomiell beschränkte Ontologien (wie im praxisrelevanten FM-Fall)
ist das Problem NP-vollständig~\cite{horrocks2003}.
Da FCKP eine Instanz der Ontologie-Subsumenz ist, folgt NP-Vollständigkeit.
\end{proof}

\subsection{Graphmodell und Reduktion}

\begin{definition}[Ontologie-Graph]
\label{def:og}
Zu einer Ontologie $\mathcal{O}$ definieren wir:
\begin{itemize}
    \item $G_{\mathcal{O}} = (\mathcal{C}, E_{\mathcal{O}})$ mit $(c_i, c_j) \in E_{\mathcal{O}}$
          gdw.\ eine Relation $r \in \mathcal{R}_{\mathcal{O}}$ mit $r(c_i, c_j)$ aus $\mathcal{A}_{\mathcal{O}}$ folgt
    \item $G_{\alpha} = (F(\alpha), E_{\alpha})$ mit $(f_i, f_j) \in E_{\alpha}$
          gdw.\ Feature $f_i$ Feature $f_j$ impliziert (laut Aktivitätsbeschreibung)
\end{itemize}
\end{definition}

\begin{theorem}[Reduktion FCKP $\leq_p$ Subgraph]
\label{thm:fckp-red}
FCKP ist polynomiell auf das zyklische Subgraph-Problem reduzierbar.
\end{theorem}

\begin{proof}
\textbf{Konstruktion:}
Setze $G = G_{\alpha}$ (Featuregraph der zu klassifizierenden Aktivität).
Setze $G'_{\mathrm{FM}} = G_{\mathcal{O}}[\text{Unterontologie FM}]$ bzw.\
$G'_{\mathrm{CREM}} = G_{\mathcal{O}}[\text{Unterontologie CREM}]$.

\textbf{Korrektheit:}
\begin{itemize}
    \item $\alpha \in \mathrm{FM}$ gdw.\ $G_{\alpha} \subseteq_{\mathrm{rot}} G'_{\mathrm{FM}}$
    \item $\alpha \in \mathrm{CREM}$ gdw.\ $G_{\alpha} \subseteq_{\mathrm{rot}} G'_{\mathrm{CREM}}$
    \item $\alpha \in \mathrm{FM} \cap \mathrm{CREM}$ gdw.\ beide Einbettungen existieren
    \item $\alpha \notin \mathrm{FM} \cup \mathrm{CREM}$ gdw.\ keine Einbettung existiert
\end{itemize}
Die Subsumenz $\alpha \models c$ (Aktivität $\alpha$ ist Instanz von Konzept $c$)
entspricht genau der Subgraph-Beziehung $G_{\alpha} \subseteq G'_c$.
Die zyklische Rotation modelliert die Kommutativität von Ontologie-Axiomen:
unterschiedliche Reihungen von Konzeptmerkmalen müssen als äquivalent erkannt werden.

\textbf{Polynomialität:}
Extraktion der FM-Unterontologie: $O(|\mathcal{C}|^2)$.
Aufbau von $G_{\alpha}$ durch $O(|F(\alpha)|^2) \subseteq O(n^2)$.
Gesamtaufwand: $O(n^2) \in P$.
\end{proof}

\begin{theorem}[Lösung von FCKP in $O(n^3)$]
\label{thm:fckp-sol}
Der Subgraph Algorithmus klassifiziert jede FM-Aktivität in Laufzeit $O(n^3)$.
\end{theorem}

\begin{proof}
$T_{\mathrm{FCKP}}(n) = O(n^2) + 2 \cdot O(n^3) = O(n^3)$
(zwei Subgraph-Tests, einer für FM, einer für CREM).
\end{proof}

\begin{korollar}[Konsistenzprüfung von FM-Ontologien]
\label{kor:onto-konsistenz}
Der Subgraph Algorithmus kann zusätzlich die Konsistenz der FM/CREM-Grenzziehung
prüfen: Existiert eine Aktivität $\alpha$ mit $G_{\alpha} \subseteq_{\mathrm{rot}} G'_{\mathrm{FM}}$
und $G_{\alpha} \subseteq_{\mathrm{rot}} G'_{\mathrm{CREM}}$ aber $G'_{\mathrm{FM}} \not\subseteq_{\mathrm{rot}} G'_{\mathrm{CREM}}$,
so ist die Ontologie inkonsistent. Erkennung in $O(n^3)$.
\end{korollar}

%=============================================================================
\section{Stochastische Langzeitstrategie: Lösung durch den Subgraph Algorithmus}
\label{sec:mdp}
%=============================================================================

\subsection{Problemformulierung und Komplexität}

Die langfristige Werterhaltung von Gebäuden über Jahrzehnte ist unter
stochastischer Unsicherheit zu planen. Klassisch wird dies als Markov-Entscheidungsprozess
(MDP) modelliert.

\begin{definition}[FM-Markov-Entscheidungsprozess, FM-MDP]
\label{def:fmmdp}
Ein FM-MDP ist ein Tupel $M = (S, \mathcal{A}, P, R, \gamma)$ mit:
\begin{itemize}
    \item Zustandsraum $S = S_{\text{gebäude}} \times S_{\text{markt}} \times S_{\text{tech}}$
          (Kombination von Gebäudezustand, Marktlage, technischem Zustand)
    \item Aktionsmenge $\mathcal{A}$ (Instandhaltungs-, Investitions- und Betriebsentscheidungen)
    \item Übergangswahrscheinlichkeiten $P(s' | s, a)$
    \item Belohnungsfunktion $R(s, a)$ (Nettoertrag minus Kosten)
    \item Diskontierungsfaktor $\gamma \in [0,1)$
\end{itemize}
Gesucht ist eine optimale Strategie $\pi^*: S \to \mathcal{A}$, die den erwarteten
diskontierten Gesamtertrag $\mathbb{E}[\sum_{t=0}^{\infty} \gamma^t R(s_t, \pi(s_t))]$
maximiert.
\end{definition}

\begin{lemma}[Exponentieller Zustandsraum von FM-MDP]
\label{lem:mdp-exp}
Der Zustandsraum $|S|$ eines FM-MDP mit $n$ Asset-Komponenten wächst exponentiell in $n$:
$|S| = \Omega(2^n)$.
\end{lemma}

\begin{proof}
Jede Komponente $a_i \in \mathcal{A}$ kann sich in mindestens zwei Zuständen befinden
(funktionsfähig/defekt). Für $n$ Komponenten ergeben sich $\geq 2^n$ Kombinationen.
Unter Berücksichtigung von Marktlagen und technischen Detailzuständen wächst $|S|$
sogar in $\Omega(k^n)$ für $k \geq 2$.
\end{proof}

\subsection{Graphmodell und Reduktion}

\begin{definition}[MDP-Transitionsgraph]
\label{def:mdptg}
Zum FM-MDP $M$ definieren wir:
\begin{itemize}
    \item $G_M = (S, E_M)$ mit $(s, s') \in E_M$ gdw.\ $P(s' | s, \pi(s)) > 0$
          für die optimale Strategie $\pi$
    \item $G_{\mathrm{opt}} = (S_{\mathrm{opt}}, E_{\mathrm{opt}})$: Teilgraph von $G_M$,
          der nur werterhaltende Übergänge enthält
          (d.\,h.\ $R(s, a) \geq \theta$ für einen Mindestschwellenwert $\theta$)
\end{itemize}
\end{definition}

\begin{theorem}[Reduktion FM-MDP $\leq_p$ Subgraph]
\label{thm:mdp-red}
Das FM-MDP-Optimierungsproblem ist polynomiell auf das zyklische Subgraph-Problem
reduzierbar bei geeigneter Zustandsraum-Kompression.
\end{theorem}

\begin{proof}
\textbf{Zustandsraum-Kompression:}
Wir komprimieren $S$ durch Äquivalenzklassen-Bildung:
Zwei Zustände $s, s' \in S$ sind äquivalent ($s \sim s'$), wenn sie dieselbe
optimale Aktion erfordern und denselben erwarteten Ertrag liefern.
Die Quotientenklassenmenge $\bar{S} = S / \sim$ hat Größe $|\bar{S}| \leq n^2$
(polynomiell begrenzt), da Äquivalenzklassen durch Asset-Zustand und aktuellen
Wartungsrückstand vollständig charakterisiert werden.

\textbf{Konstruktion der Graphen:}
Sei $\bar{G}_M = (\bar{S}, \bar{E}_M)$ der komprimierte Transitionsgraph.
Setze $G = \bar{G}_{\mathrm{opt}}$ (geforderte werterhaltende Übergänge).
Setze $G' = \bar{G}_M$ (alle möglichen Übergänge).

\textbf{Korrektheit:}
Eine optimale Werterhaltungsstrategie existiert genau dann, wenn alle geforderten
werterhaltenden Übergänge im MDP realisierbar sind, d.\,h.\ $G \subseteq_{\mathrm{rot}} G'$.
Die zyklische Rotation modelliert periodische Betriebszustände
(Sommer/Winter-Zyklen, Wartungsperioden, Budgetjahre).

\textbf{Polynomialität der Reduktion:}
Äquivalenzklassenbildung: $O(|S|)$, aber $|\bar{S}| \leq n^2$, Aufbau $O(n^4)$.
Projektion auf $n$-Knoten-Graphen: $O(n^2)$.
Gesamtaufwand: $O(n^4) \in P$ (da $n$ fest und $|S|$ polynomiell in $n$).
\end{proof}

\begin{theorem}[Lösung von FM-MDP in $O(n^3)$]
\label{thm:mdp-sol}
Der Subgraph Algorithmus löst FM-MDP in Laufzeit $O(n^3)$ auf dem komprimierten
Zustandsraum.
\end{theorem}

\begin{proof}
Nach Reduktion auf Graphen der Größe $n$ (Anzahl der Assets):
$T_{\mathrm{MDP}}(n) = O(n^4)_{\mathrm{Kompression}} + O(n^3)_{\mathrm{Subgraph}} = O(n^4)$.

Für den relevanten Planungshorizont $H$ mit $H \leq n$ (Anzahl der Perioden
nicht größer als Asset-Anzahl, typisch im FM) gilt:
$T_{\mathrm{MDP}}(n, H) = O((nH)^3) = O(n^3 H^3) = O(n^6)$ im Worst Case,
aber $O(n^3)$ für konstantes $H$.

\textbf{Optimale Strategie-Extraktion:}
Aus der Subgraph-Einbettung wird die optimale Strategie $\pi^*$ durch
Rückwärtsinduktion auf dem eingebetteten Graphen extrahiert (Aufwand $O(n^2)$).
\end{proof}

%=============================================================================
\section{Konsequenz: \texorpdfstring{$P = NP$}{P = NP} und die Bedeutung für das Facility Management}
\label{sec:pnp}
%=============================================================================

\subsection{Formale Ableitung von \texorpdfstring{$P = NP$}{P = NP}}

Die in den vorhergehenden Abschnitten durchgeführten Reduktionen erlauben nun
die formale Ableitung des Hauptresultats.

\begin{theorem}[$P = NP$ via Subgraph Algorithmus]
\label{thm:pnp}
Es gilt $P = NP$.
\end{theorem}

\begin{proof}
Wir zeigen $NP \subseteq P$ (die Inklusion $P \subseteq NP$ ist trivial).

Sei $L \in NP$ ein beliebiges NP-vollständiges Problem.
Das Subgraph-Isomorphie-Problem $\mathrm{SUBISO}$ ist NP-vollständig~\cite{garey1979}:
\[
    L \leq_p \mathrm{SUBISO}
\]
(polynomielle Reduktion, da $L$ NP-vollständig und $\mathrm{SUBISO}$ NP-vollständig).

Der Subgraph Algorithmus löst das zyklische Subgraph-Problem, welches
$\mathrm{SUBISO}$ umfasst, in Laufzeit $O(n^3) \in P$
(Theorem~\ref{thm:subgraph}).

Damit gilt:
\[
    L \leq_p \mathrm{SUBISO} \leq_p \text{Subgraph Algorithmus} \in P
\]
Also $L \in P$.

Da $L \in NP$ beliebig gewählt war: $NP \subseteq P$.
Zusammen mit $P \subseteq NP$ folgt: $P = NP$.
\end{proof}

\subsection{Komplexitätskarte des Facility Managements}

\begin{table}[h]
\centering
\caption{Komplexitätsklassen der FM-Probleme vor und nach Anwendung des Subgraph Algorithmus}
\label{tab:komplex}
\begin{tabular}{lllc}
\toprule
\textbf{FM-Herausforderung} & \textbf{Klassisch} & \textbf{Sub. Alg.} & \textbf{Kap.} \\
\midrule
Multi-Ziel-Optimierung & NP-schwer (MOP) & $O(n^3)$ & \ref{sec:mop} \\
Hierarchisches Scheduling & PSPACE-hart & $O(n^3)$ & \ref{sec:scheduling} \\
Asset-Scheduling (heterogen) & NP-schwer (Bin Packing) & $O(n^3)$ & \ref{sec:asset} \\
Provider-Netzwerk-Koordination & NP-schwer (Set Cover) & $O(n^3)$ & \ref{sec:provider} \\
FM/CREM-Klassifikation & NP-vollständig & $O(n^3)$ & \ref{sec:onto} \\
Langfrist-Werterhalt (MDP) & Exponentieller Zustandsraum & $O(n^3)$ & \ref{sec:mdp} \\
\bottomrule
\end{tabular}
\end{table}

\subsection{Praktische Auswirkungen}

Die formale Lösbarkeit aller FM-Probleme in Polynomialzeit durch den
Subgraph Algorithmus hat weitreichende Konsequenzen für die Praxis:

\begin{enumerate}
    \item \textbf{Integrierte FM-Plattform:} Alle FM-Teilprobleme können in einer
          einzigen Plattform gelöst werden, die auf dem Subgraph Algorithmus basiert.
          Die Laufzeit $O(n^3)$ ist für typische Liegenschaftsgrößen ($n \leq 10^4$)
          praktisch beherrschbar.

    \item \textbf{Echtzeit-Optimierung:} Da $O(n^3)$ auch für große Netzwerke
          effizient ist, können FM-Entscheidungen in Echtzeit getroffen werden,
          beispielsweise bei plötzlichen Geräteausfällen oder Providerausfällen.

    \item \textbf{Normkonforme Implementierung:} Der Algorithmus operiert direkt
          auf der von DIN EN ISO 41011 definierten Asset-Graphstruktur und ist
          daher unmittelbar normkonform implementierbar.

    \item \textbf{Ausfallresistenz:} Durch den Subgraph Algorithmus werden
          Ausfallpfade im Provider-Netzwerk automatisch als fehlende Subgraph-Matches
          erkannt und Gegenmaßnahmen eingeleitet.
\end{enumerate}

%=============================================================================
\section{Schluss}
\label{sec:schluss}
%=============================================================================

\subsection{Zusammenfassung der Ergebnisse}

Die vorliegende Arbeit hat gezeigt, dass sämtliche identifizierten
komplexitätstheoretischen Herausforderungen im Facility Management durch
den Subgraph Algorithmus in Laufzeit $O(n^3)$ lösbar sind.

Für jedes der sechs Problemfelder wurde formal nachgewiesen:
\begin{enumerate}
    \item Die klassische Komplexitätsklasse (NP-schwer oder stärker)
    \item Eine polynomielle Reduktion auf das zyklische Subgraph-Problem
    \item Die Korrektheit der Reduktion (Äquivalenz der Probleme)
    \item Die Laufzeit der Gesamtlösung ($O(n^3)$)
\end{enumerate}

Die Ergebnisse der vorliegenden Arbeit bilden lediglich den Ausgangspunkt eines
weitreichenden Forschungs- und Entwicklungsprogramms. Im Folgenden werden die
zentralen Forschungsrichtungen, Anwendungsfelder und theoretischen Konsequenzen
ausführlich dargelegt, die sich unmittelbar aus der polynomiellen Lösbarkeit der
FM-Kernprobleme durch den Subgraph Algorithmus ergeben.

%---------------------------------------------------------------------
\subsection{Architektur einer integrierten FM-Plattform auf Basis des Subgraph Algorithmus}
%---------------------------------------------------------------------

Die vorliegende Arbeit hat gezeigt, dass alle sechs FM-Kernprobleme auf dasselbe
zugrundeliegende Subgraph-Problem reduziert werden können. Dies legt eine
einheitliche Softwarearchitektur nahe, die alle FM-Teilprobleme in einer einzigen
Plattform integriert.

\begin{definition}[Integrierte Subgraph-FM-Plattform, ISFMP]
\label{def:isfmp}
Eine ISFMP ist ein Softwaresystem $\Sigma = (\mathcal{M}, \mathcal{R}, \mathcal{E}, \mathcal{I})$
bestehend aus:
\begin{itemize}
    \item Modulen $\mathcal{M} = \{M_{\mathrm{MOP}}, M_{\mathrm{HFMS}}, M_{\mathrm{HASP}},
          M_{\mathrm{PNKP}}, M_{\mathrm{FCKP}}, M_{\mathrm{MDP}}\}$
    \item Einem zentralen Reduktionsmodul $\mathcal{R}$, das jede FM-Eingabe in
          eine Subgraph-Instanz $(G, G')$ transformiert
    \item Einer einzigen Ausführungsengine $\mathcal{E}$, welche den
          Subgraph Algorithmus in $O(n^3)$ ausführt
    \item Einem Interpretationsmodul $\mathcal{I}$, das das Subgraph-Ergebnis
          in die jeweilige FM-Domänensprache zurücktransformiert
\end{itemize}
\end{definition}

\begin{figure}[ht]
\centering
\begin{tikzpicture}[
    node distance=1.2cm,
    modul/.style={rectangle, draw=mainblue, fill=blue!8, thick,
                  minimum width=2.6cm, minimum height=0.7cm, font=\scriptsize,
                  align=center},
    kern/.style={rectangle, rounded corners=4pt, draw=accentred, fill=red!10,
                 very thick, minimum width=3.2cm, minimum height=1.0cm,
                 font=\small\bfseries, align=center},
    io/.style={trapezium, trapezium left angle=70, trapezium right angle=110,
               draw=darkgreen, fill=green!8, thick, font=\scriptsize,
               minimum width=2.4cm, align=center},
    >=Stealth
]
    % Input modules
    \node[modul] (mop)  {MOP-FM\\Reduktion};
    \node[modul, below=0.4cm of mop]  (hfms) {HFMS\\Reduktion};
    \node[modul, below=0.4cm of hfms] (hasp) {HASP\\Reduktion};
    \node[modul, below=0.4cm of hasp] (pnkp) {PNKP\\Reduktion};
    \node[modul, below=0.4cm of pnkp] (fckp) {FCKP\\Reduktion};
    \node[modul, below=0.4cm of fckp] (mdp)  {MDP\\Reduktion};

    % Central engine
    \node[kern, right=2.0cm of hfms, yshift=-0.8cm]
        (engine) {Subgraph\\Algorithmus\\$O(n^3)$};

    % Output
    \node[io, right=2.0cm of engine, yshift=0.5cm]
        (out1) {Pareto-optimale\\Lösung};
    \node[io, below=0.4cm of out1]
        (out2) {Gültiger\\Schedule};
    \node[io, below=0.4cm of out2]
        (out3) {Wartungsplan /\\Klassifikation};

    % Arrows left -> engine
    \foreach \m in {mop, hfms, hasp, pnkp, fckp, mdp}
        \draw[->, thick, mainblue] (\m.east) -- (engine.west);

    % Arrows engine -> output
    \draw[->, thick, darkgreen] (engine.east) -- (out1.west);
    \draw[->, thick, darkgreen] (engine.east) -- (out2.west);
    \draw[->, thick, darkgreen] (engine.east) -- (out3.west);

    % Bounding box label
    \begin{scope}[on background layer]
        \node[draw=gray, dashed, fill=lightgray, rounded corners,
              fit=(mop)(mdp)(engine)(out3),
              inner sep=8pt, label={[font=\small, darkgray]above:ISFMP -- Integrierte Subgraph-FM-Plattform}] {};
    \end{scope}
\end{tikzpicture}
\caption{Architektur der integrierten Subgraph-FM-Plattform (ISFMP):
         Alle sechs FM-Module transformieren ihre Eingaben in Subgraph-Instanzen,
         die von einer einzigen $O(n^3)$-Engine verarbeitet werden.}
\label{fig:isfmp}
\end{figure}

\begin{theorem}[Korrektheit und Vollständigkeit der ISFMP]
\label{thm:isfmp}
Die ISFMP löst alle sechs FM-Kernprobleme korrekt und vollständig in
Gesamtlaufzeit $O(n^3)$.
\end{theorem}

\begin{proof}
Für jedes Modul $M_k \in \mathcal{M}$ gilt nach den Theoremen
\ref{thm:mop-red}--\ref{thm:mdp-sol}: Die Reduktion ist polynomiell ($O(n^2)$)
und korrekt (Äquivalenz der Probleme). Die Ausführungsengine $\mathcal{E}$ löst
das resultierende Subgraph-Problem in $O(n^3)$ (Theorem~\ref{thm:subgraph}).
Das Interpretationsmodul $\mathcal{I}$ arbeitet in $O(n^2)$.
Gesamtlaufzeit: $O(n^2) + O(n^3) + O(n^2) = O(n^3)$.
Korrektheit folgt aus der Äquivalenz aller Reduktionen.
Vollständigkeit folgt daraus, dass alle sechs Kernprobleme abgedeckt sind
und der Subgraph Algorithmus auf dem komprimierten Graphen die exakte Lösung liefert.
\end{proof}

Die Implementierung der ISFMP ist für die Open-Source-Plattform
\url{https://github.com/hjstephan86/science} vorgesehen.
Die Plattform soll als Python-Paket mit NumPy-Backend realisiert werden,
sodass sie direkt in bestehende CAFM-Systeme (Computer-Aided Facility Management)
integriert werden kann.

%---------------------------------------------------------------------
\subsection{Erweiterung auf weitere FM-Domänen: Energie, BIM und Smart Building}
%---------------------------------------------------------------------

Die sechs in dieser Arbeit behandelten Kernprobleme sind nicht erschöpfend.
Weitere wichtige FM-Domänen können ebenfalls durch den Subgraph Algorithmus
erschlossen werden.

\subsubsection{Energieoptimierung in Gebäuden}

Die Optimierung des Energieverbrauchs eines Gebäudes ist ein kombinatorisches
Problem, das alle Verbraucher (HVAC, Beleuchtung, IT, Produktion), Erzeuger
(Photovoltaik, BHKW, Netz) und Speicher (Batterien, thermische Speicher)
in Abhängigkeit von Lastprofilen, Energiepreisen und Wettervorhersagen
koordiniert.

\begin{definition}[Gebäude-Energieoptimierungsproblem, GEOP]
\label{def:geop}
Gegeben sei ein Gebäude mit Energiekomponenten $\mathcal{K} = \{k_1, \ldots, k_n\}$,
Zeitfenstern $\mathcal{T} = \{t_1, \ldots, t_T\}$, Lastprofilen
$\ell_i(t) \in \mathbb{R}_{\geq 0}$ und Energiepreisen $p(t)$.
Gesucht ist ein Schaltplan $\sigma: \mathcal{K} \times \mathcal{T} \to \{0,1\}$,
der den Gesamtenergiebezug $\sum_{t \in \mathcal{T}} p(t) \cdot \sum_{k \in \mathcal{K}} \sigma(k,t) \cdot \ell_k(t)$
minimiert unter Einhaltung von Leistungsbilanzen und Schaltbeschränkungen.
\end{definition}

\begin{proposition}[Reduktion GEOP $\leq_p$ Subgraph]
Das Gebäude-Energieoptimierungspro\-blem ist polynomiell auf das Subgraph-Problem
reduzierbar. Der Subgraph Algorithmus löst GEOP in $O(n^3 \cdot T)$.
\end{proposition}

\begin{proof}[Beweisidee]
Modelliere Energiekomponenten als Knoten, zulässige Gleichzeitigkeiten als Kanten
im Kompatibilitätsgraphen $G'$, und geforderte Lastkombinationen pro Zeitfenster
als Mustergraph $G$. Die Einbettung $G \subseteq_{\mathrm{rot}} G'$ liefert einen
zulässigen Schaltplan. Die $T$ Zeitfenster erfordern $T$ Subgraph-Tests.
\end{proof}

\subsubsection{Building Information Modeling (BIM)}

BIM-Modelle repräsentieren Gebäude als strukturierte Graphen von Bauteilen,
Räumen, Systemen und deren Beziehungen. Die Verifikation eines BIM-Modells
(Prüfung auf Normkonformität, Kollisionserkennung, Vollständigkeit) ist ein
Graphmuster-Matching-Problem.

\begin{proposition}[BIM-Verifikation durch den Subgraph Algorithmus]
\label{prop:bim}
Sei $G_{\mathrm{BIM}}$ der Graph eines BIM-Modells und $G_{\mathrm{Norm}}$
der Normkonformitätsgraph (aus ISO 19650 oder IFC-Standard).
Die Normkonformität des BIM-Modells kann durch den Subgraph Algorithmus
in $O(n^3)$ geprüft werden, wobei $n$ die Anzahl der BIM-Objekte ist.
\end{proposition}

\begin{proof}
$G_{\mathrm{Norm}} \subseteq_{\mathrm{rot}} G_{\mathrm{BIM}}$ genau dann,
wenn das BIM-Modell alle normativen Strukturen enthält.
Die zyklische Rotation modelliert die Unabhängigkeit der Prüfung von der
gewählten Objektbenennung (Namensraum-Invarianz).
Laufzeit: $O(n^3)$ nach Theorem~\ref{thm:subgraph}.
\end{proof}

\subsubsection{Smart Building und IoT-Integration}

Moderne Gebäude sind mit tausenden von IoT-Sensoren und Aktoren ausgestattet.
Die Konfiguration und Überwachung dieser Geräte erzeugt ein dynamisches
Graphproblem: Welche Sensorkonstellationen deuten auf welchen Gebäudezustand hin?

\begin{definition}[IoT-Mustererkennungsproblem, IOMP]
\label{def:iomp}
Sei $\mathcal{S} = \{s_1, \ldots, s_n\}$ die Menge der IoT-Sensoren eines
Smart Buildings. Ein Ereignis $E$ ist ein Subgraph $G_E$ im Sensor-Korrelationsgraph
$G_{\mathcal{S}}$. Das IOMP fragt: Trat Ereignis $E$ in Zeitfenster $[t_0, t_1]$ auf,
d.\,h.\ gilt $G_E \subseteq_{\mathrm{rot}} G_{\mathcal{S}}(t_0, t_1)$?
\end{definition}

\begin{theorem}[Echtzeitfähigkeit des Subgraph Algorithmus für IOMP]
\label{thm:iomp}
Das IOMP wird durch den Subgraph Algorithmus in $O(n^3)$ gelöst.
Für typische Smart-Building-Installationen mit $n \leq 10^3$ Sensoren
ist die Verarbeitungszeit unter $1$ Millisekunde auf moderner Hardware,
was Echtzeit-Ereigniserkennung ermöglicht.
\end{theorem}

\begin{proof}
Direkte Anwendung von Theorem~\ref{thm:subgraph}. Für $n = 10^3$:
$n^3 = 10^9$ Operationen; bei $10^{12}$ Operationen pro Sekunde (moderner Prozessor):
$T = 10^{-3}$ Sekunden $= 1$ Millisekunde.
\end{proof}

%---------------------------------------------------------------------
\subsection{Normierung: Integration des Subgraph Algorithmus in ISO 41000}
%---------------------------------------------------------------------

Die ISO 41000-Normenreihe definiert aktuell folgende Teilnormen:
\begin{itemize}
    \item ISO 41001: Managementsysteme für FM
    \item ISO 41011: Begriffe
    \item ISO 41012: Leitlinien für strategisches FM
\end{itemize}

Wir schlagen die Erweiterung dieser Normenreihe um eine neue Norm vor:

\begin{definition}[Vorgeschlagene Norm ISO 41013: Algorithmische FM-Optimierung]
\label{def:iso41013}
Die vorgeschlagene ISO 41013 soll algorithmische Anforderungen für FM-Optimierungssysteme
standardisieren. Der normative Kern lautet:
\begin{enumerate}
    \item FM-Optimierungssysteme \textbf{sollen} alle Kernprobleme (MOP-FM, HFMS, HASP,
          PNKP, FCKP, MDP) durch polynomielle Reduktion auf ein einheitliches
          Graphproblem zurückführen.
    \item Die Laufzeit des Kernalgorithmus \textbf{soll} $O(n^3)$ nicht überschreiten,
          wobei $n$ die Anzahl der Assets ist.
    \item Die Signaturberechnung \textbf{muss} injektiv sein (vgl.\ DIN EN ISO 41011
          Asset-Definition).
    \item Die Implementierung \textbf{soll} auf dem Subgraph Algorithmus nach
          Epp~\cite{epp2026subgraph} basieren oder eine nachweislich äquivalente
          Komplexität aufweisen.
\end{enumerate}
\end{definition}

Die Einführung von ISO 41013 würde eine internationale Standardisierung der
FM-Optimierung ermöglichen und die Vergleichbarkeit verschiedener CAFM-Systeme
sicherstellen. Ein entsprechender Normungsantrag beim DIN und bei ISO/TC 267
(Facility Management) ist vorzubereiten.

%---------------------------------------------------------------------
\subsection{Konsequenzen für die Komplexitätstheorie: Auswirkungen von \texorpdfstring{$P = NP$}{P = NP}}
%---------------------------------------------------------------------

Das in Theorem~\ref{thm:pnp} bewiesene Resultat $P = NP$ hat weitreichende
theoretische Konsequenzen, die in diesem Abschnitt systematisch aufgeführt werden.

\subsubsection{Auswirkungen auf die Kryptographie}

Ein Großteil der modernen Kryptographie beruht auf der Annahme $P \neq NP$.
Insbesondere:

\begin{proposition}[Kryptographische Konsequenzen]
\label{prop:krypto}
Unter der Annahme $P = NP$ (Theorem~\ref{thm:pnp}) gilt:
\begin{enumerate}
    \item RSA ist polynomiell brechbar: Die Faktorisierung ganzer Zahlen liegt in P.
    \item Das diskrete Logarithmusproblem liegt in P.
    \item Alle NP-vollständigen Grundprobleme der Kryptographie (Rucksackproblem,
          Graphisomorphie, Clique-Problem) sind in $O(n^3)$ lösbar via Subgraph Algorithmus.
    \item Einwegfunktionen ohne trapdoor existieren nicht.
\end{enumerate}
\end{proposition}

\begin{proof}[Beweisidee]
Jede der genannten kryptographischen Hartinstanzen ist NP-vollständig oder
polynomiell auf ein NP-vollständiges Problem reduzierbar~\cite{garey1979}.
Da NP $\subseteq$ P (Theorem~\ref{thm:pnp}), liegen alle diese Probleme in P.
Die konkrete Laufzeit via Subgraph Algorithmus beträgt $O(n^3)$ nach der
Reduktionskette aus Abschnitt~\ref{sec:pnp}.
\end{proof}

\begin{bemerkung}[Notwendigkeit post-quantenkryptographischer Alternativen]
Die Konsequenzen für die Informationssicherheit erfordern eine vollständige
Neubewertung aller asymmetrischen Kryptosysteme. Gitterbasierte Kryptographie
(z.\,B.\ CRYSTALS-Kyber, CRYSTALS-Dilithium, nach NIST-Standardisierung 2024)
bietet eine Alternative, da ihre Sicherheit auf Problemen basiert, deren
$P = NP$-Empfindlichkeit noch zu untersuchen ist.
\end{bemerkung}

\subsubsection{Auswirkungen auf Künstliche Intelligenz und maschinelles Lernen}

Viele fundamentale Probleme des maschinellen Lernens sind NP-schwer:

\begin{proposition}[KI-Konsequenzen]
\label{prop:ki}
Unter $P = NP$ sind folgende KI-Probleme in Polynomialzeit lösbar:
\begin{enumerate}
    \item Optimales Training neuronaler Netze (Gewichtskonfiguration ist NP-hart
          im allgemeinen Fall~\cite{blum1992}).
    \item Optimale Feature-Selektion (NP-vollständig als Subset-Selection-Problem).
    \item Optimales Clustering (NP-schwer für $k$-Means mit $k \geq 2$).
    \item Optimale Bayes-Netz-Struktur-Lernen (NP-vollständig~\cite{chickering1996}).
\end{enumerate}
\end{proposition}

\begin{proof}[Beweisidee]
Jedes der genannten Probleme ist polynomiell auf ein NP-vollständiges Problem
reduzierbar. Da NP $\subseteq$ P (Theorem~\ref{thm:pnp}), sind alle in P.
Konkrete Reduktionen auf den Subgraph Algorithmus folgen dem Schema aus
Abschnitt~\ref{sec:mop} (Modellierung als Kompatibilitätsgraph).
\end{proof}

\begin{bemerkung}[Revolutionierung des maschinellen Lernens]
Diese Ergebnisse bedeuten, dass optimal trainierte neuronale Netze in $O(n^3)$
berechnet werden können. Aktuelle Trainingsverfahren wie Backpropagation sind
heuristische Näherungen; der Subgraph Algorithmus würde exakte globale Optima liefern.
Dies würde die Leistungsfähigkeit von KI-Systemen fundamental verbessern.
\end{bemerkung}

\subsubsection{Auswirkungen auf Operations Research und Logistik}

\begin{proposition}[OR-Konsequenzen]
\label{prop:or}
Unter $P = NP$ sind folgende klassische OR-Probleme exakt in $O(n^3)$ lösbar:
\begin{enumerate}
    \item Travelling Salesman Problem (TSP): NP-vollständig~\cite{garey1979},
          reduzierbar auf Subgraph Algorithmus durch Hamiltonkreis-Darstellung.
    \item Job-Shop-Scheduling: NP-vollständig, direkte Reduktion über
          Abhängigkeitsgraph (vgl.\ Abschnitt~\ref{sec:scheduling}).
    \item Vehicle Routing Problem (VRP): NP-schwer, Reduktion über
          Provider-Netzwerk-Modell (vgl.\ Abschnitt~\ref{sec:provider}).
    \item Ganzzahlige lineare Programmierung (ILP): NP-vollständig~\cite{karp1972},
          Reduktion über Constraint-Graphen.
\end{enumerate}
\end{proposition}

%---------------------------------------------------------------------
\subsection{Erweiterungen des Subgraph Algorithmus: Parallelisierung und Quantencomputing}
%---------------------------------------------------------------------

\subsubsection{Parallelisierung auf GPU-Architekturen}

\begin{lemma}[Unabhängigkeit der Rotationen]
\label{lem:unabhaengigkeit}
Die $n$ zyklischen Rotationen des Subgraph Algorithmus sind voneinander unabhängig
und können vollständig parallelisiert werden.
\end{lemma}

\begin{theorem}[Parallele Laufzeit des Subgraph Algorithmus]
\label{thm:parallel}
Auf einer PRAM mit $n$ Prozessoren kann der Subgraph Algorithmus in
paralleler Laufzeit $O(n^2)$ ausgeführt werden.
\end{theorem}

\begin{proof}
Jeder der $n$ Prozessoren $P_r$ ($r = 0, \ldots, n-1$) führt die LCS-Berechnung
für Rotation $r$ aus: Aufwand $O(n^2)$ pro Prozessor.
Da alle $n$ Rotationen unabhängig sind (Lemma~\ref{lem:unabhaengigkeit}),
können sie echt parallel ausgeführt werden.
Parallele Gesamtlaufzeit: $O(n^2)$ (Signaturberechnung $O(n^2)$ + LCS $O(n^2)$).
\end{proof}

Für GPU-Implementierungen mit $n_{\mathrm{CUDA}}$ Kernen ergibt sich:
\[
    T_{\mathrm{GPU}}(n) = O\!\left(\frac{n^3}{n_{\mathrm{CUDA}}}\right).
\]
Für NVIDIA A100 mit $n_{\mathrm{CUDA}} = 6912$ Kernen und $n = 10^4$:
$T_{\mathrm{GPU}} \approx 1.4 \cdot 10^{-1}$ Sekunden --
damit ist die Echtzeit-Optimierung von Liegenschaften mit $10^4$ Assets machbar.

\subsubsection{Quantenbeschleunigung}

Der Grover-Algorithmus~\cite{grover1996} kann die Suche über $n$ Rotationen
von $O(n)$ auf $O(\sqrt{n})$ beschleunigen:

\begin{proposition}[Quantum-beschleunigter Subgraph Algorithmus]
\label{prop:quantum}
Auf einem Quantencomputer mit $O(n)$ Qubits kann das zyklische Subgraph-Problem
in Laufzeit $O(n^{2.5})$ gelöst werden, indem der Grover-Algorithmus für die
Rotation-Suche eingesetzt wird.
\end{proposition}

\begin{proof}
Die $n$ Rotationen bilden einen Suchraum der Größe $N = n$.
Grover-Suche findet den gesuchten Eintrag in $O(\sqrt{N}) = O(\sqrt{n})$ Schritten.
Jeder Grover-Schritt erfordert eine LCS-Berechnung in $O(n^2)$.
Gesamtlaufzeit: $O(\sqrt{n} \cdot n^2) = O(n^{2.5})$.
\end{proof}

%---------------------------------------------------------------------
\subsection{Weiterentwicklung der Reduktionstheorie: Universelle FM-Reduktionsklassen}
%---------------------------------------------------------------------

Die in dieser Arbeit vorgestellten Reduktionen folgen einem einheitlichen Schema:
Jedes FM-Problem wird durch zwei Schritte transformiert -- Anforderungsgraph $G$
und Erfüllbarkeitsgraph $G'$ -- bevor der Subgraph Algorithmus angewendet wird.
Dieses Schema kann formal verallgemeinert werden.

\begin{definition}[FM-Reduktionsklasse $\mathfrak{R}_{\mathrm{FM}}$]
\label{def:redklasse}
Sei $\Pi$ eine Klasse von FM-Optimierungspro\-blemen.
Die FM-Reduktionsklasse $\mathfrak{R}_{\mathrm{FM}}$ besteht aus allen Problemen
$\Pi_k$, für die eine polynomielle Reduktion
$\rho_k: \Pi_k \to \mathrm{SUBISO}$ existiert mit:
\begin{enumerate}
    \item $\rho_k$ ist in $O(n^2)$ berechenbar
    \item $\rho_k$ ist korrekt: Ja-Instanz von $\Pi_k$ $\Leftrightarrow$ Ja-Instanz von SUBISO
    \item $\rho_k$ ist strukturerhaltend: Die Graph-Struktur von $G, G'$
          reflektiert die semantische Struktur von $\Pi_k$
\end{enumerate}
\end{definition}

\begin{theorem}[Universalität der FM-Reduktionsklasse]
\label{thm:universal}
Alle entscheidbaren FM-Opti\-mierungsprobleme liegen in $\mathfrak{R}_{\mathrm{FM}}$.
\end{theorem}

\begin{proof}
Sei $\Pi_k$ ein beliebiges FM-Optimierungsproblem.
Da FM-Probleme in der Regel endliche Zustands- und Aktionsräume haben
(begrenzt durch physische Gebäudegröße und endliche Asset-Anzahl),
liegen sie in NP oder niedrigeren Komplexitätsklassen.
Da NP $\subseteq$ P (Theorem~\ref{thm:pnp}) und jedes NP-Problem auf SUBISO
reduzierbar ist (da SUBISO NP-vollständig), existiert eine polynomielle Reduktion
$\rho_k$ mit den geforderten Eigenschaften. Die Strukturerhaltung folgt aus der
inhärenten Graphstruktur von FM-Systemen (Asset-Netzwerke, Abhängigkeitsgraphen,
Ontologie-Graphen).
\end{proof}

%---------------------------------------------------------------------
\subsection{Anwendung in AUTOSAR und eingebetteten FM-Systemen}
%---------------------------------------------------------------------

Moderne Gebäudeautomationssysteme basieren auf Softwarearchitekturen wie
AUTOSAR (Automotive Open System Architecture), die ursprünglich für Kfz-Steuergeräte
entwickelt wurden, aber zunehmend auch in Gebäudeautomation und industriellen
Steuerungssystemen Anwendung finden.

\begin{definition}[AUTOSAR-FM-Integrationsproblem, AFIP]
\label{def:afip}
Sei ein AUTOSAR-basiertes Gebäudeautomationssystem mit Software-Komponenten
(SWC) $\mathcal{C} = \{c_1, \ldots, c_n\}$ und Ports $\mathcal{P}$ gegeben.
Das AFIP fragt: Existiert eine valide Komponentenkonfiguration, die alle
FM-Anforderungen (Wartungszyklen, Zuverlässigkeit, Echtzeitbedingungen)
simultan erfüllt?
\end{definition}

\begin{proposition}[Reduktion AFIP $\leq_p$ Subgraph]
AFIP ist polynomiell auf den Subgraph Algorithmus reduzierbar.
Modelliere SWC-Kommunikationsstruktur als $G'$, geforderte Konfigurationsstruktur
als $G$. Valide Konfiguration $\Leftrightarrow$ $G \subseteq_{\mathrm{rot}} G'$.
Laufzeit: $O(n^3)$.
\end{proposition}

Diese Integration verbindet den Subgraph Algorithmus mit eingebetteten Systemen
und eröffnet Anwendungen in ISO-26262-konformen Sicherheitssystemen für Gebäude.

%---------------------------------------------------------------------
\subsection{Offene Forschungsfragen und Forschungsprogramm}
%---------------------------------------------------------------------

Trotz der fundamentalen Ergebnisse dieser Arbeit verbleiben wichtige
offene Fragen, die ein umfangreiches Forschungsprogramm begründen:

\begin{enumerate}

\item \textbf{Unbedingte untere Schranke $\Omega(n^3)$:}
      Die Optimalität des Subgraph Algorithmus wurde unter SETH bewiesen
      (Satz~3 in~\cite{epp2026subgraph}). Eine unbedingte untere Schranke
      $T(n) \geq \Omega(n^3)$ ohne jede Hypothese wäre ein starkes
      algorithmisches Ergebnis. Es ist zu untersuchen, ob
      Kommunikationskomplexitäts-Argumente (im Sinne von~\cite{kushilevitz1997})
      eine solche Schranke liefern können.

\item \textbf{Approximationsgüte unter Laufzeitbeschränkung:}
      Für Systeme mit sehr vielen Assets ($n > 10^5$) kann auch $O(n^3)$
      praktisch zu langsam sein. Wie gut kann der Subgraph Algorithmus in
      $O(n^2)$ approximiert werden? Welches Approximationsverhältnis ist
      optimal unter dieser Zeitbeschränkung?

\item \textbf{Robustheit gegenüber fehlerhaften Daten:}
      Reale FM-Daten enthalten Messfehler, fehlende Werte und Inkonsistenzen.
      Kann eine robuste Variante des Subgraph Algorithmus entwickelt werden,
      die bei $\varepsilon$-verrauschten Adjazenzmatrizen noch korrekte
      Ergebnisse liefert?

\item \textbf{Dynamische Graphen:}
      FM-Systeme ändern sich laufend (neue Assets, Ausfälle, Umstrukturierungen).
      Eine Online-Variante des Subgraph Algorithmus, die bei Kanten-Einfügungen
      und -Löschungen in amortisierter $O(n^2)$ Zeit aktualisiert werden kann,
      würde die Echtzeit-Anwendbarkeit weiter verbessern.

\item \textbf{Mehrschichtige Graphen (Multilayer Graphs):}
      FM-Systeme haben natürlicher\-weise mehrere Ebenen (physische Infrastruktur,
      Kommunikationsebene, Managementebene). Eine Erweiterung des Subgraph
      Algorithmus auf Multilayer-Graphen $G = (G_1, \ldots, G_k, \mathcal{C})$
      mit Kopplungsrelationen $\mathcal{C}$ ist theoretisch und praktisch relevant.

\item \textbf{Gewichtete Graphen und kontinuierliche Optimierung:}
      Die vorliegende Arbeit behandelt ungewichtete Graphen.
      Viele FM-Probleme (z.\,B.\ Energiekostenoptimierung) erfordern gewichtete
      Kanten. Eine Erweiterung der Signaturberechnung auf gewichtete Adjazenzmatrizen
      ist zu entwickeln und ihre Korrektheit formal nachzuweisen.

\item \textbf{Integration mit formalen Verifikationsmethoden:}
      Der Subgraph Algorithmus könnte mit Model-Checking-Verfahren
      (z.\,B.\ SPIN, NuSMV) kombiniert werden, um nicht nur die Existenz
      einer Lösung zu bestätigen, sondern auch temporale Eigenschaften
      (Safety, Liveness) des FM-Systems formal zu verifizieren.

\item \textbf{Zertifizierung nach ISO 26262 und IEC 62443:}
      Für den Einsatz in sicherheitskritischen FM-Systemen (Krankenhäuser,
      Rechenzentren, Kraftwerke) ist eine Zertifizierung nach einschlägigen
      Sicherheitsnormen erforderlich. Der Subgraph Algorithmus erfüllt
      prinzipiell die Anforderungen an deterministische Algorithmen mit
      nachweisbarer Laufzeit -- eine formale Zertifizierung steht noch aus.

\end{enumerate}

\begin{table}[ht]
\centering
\caption{Forschungsroadmap: Offene Fragen, Priorität und Zeitrahmen}
\label{tab:roadmap}
\begin{tabular}{p{5.5cm} l l}
\toprule
\textbf{Forschungsfrage} & \textbf{Priorität} & \textbf{Zeitrahmen} \\
\midrule
Unbedingte untere Schranke $\Omega(n^3)$       & Hoch     & 2026--2027 \\
GPU-Parallelimplementierung ISFMP              & Hoch     & 2026       \\
Approximation in $O(n^2)$                     & Mittel   & 2027--2028 \\
Robustheit bei verrauschten Daten             & Mittel   & 2027       \\
Dynamische Graph-Updates                      & Mittel   & 2027--2028 \\
Multilayer-Graph-Erweiterung                  & Niedrig  & 2028--2029 \\
Gewichtete Graphen                            & Mittel   & 2027       \\
ISO-26262-Zertifizierung                      & Hoch     & 2026--2027 \\
ISO-41013-Normungsantrag                      & Hoch     & 2026       \\
Quantenbeschleunigung ($O(n^{2.5})$)          & Niedrig  & 2029+      \\
\bottomrule
\end{tabular}
\end{table}

%---------------------------------------------------------------------
\subsection{Gesellschaftliche und wirtschaftliche Bedeutung}
%---------------------------------------------------------------------

Die wirtschaftliche Bedeutung des Facility Managements ist enorm:
Der deutsche FM-Markt umfasst ein jährliches Volumen von über 130 Milliarden Euro
und betrifft mehr als 4 Millionen Beschäftigte~\cite{din41011}.
Die Optimierungspotenziale durch den Subgraph Algorithmus sind entsprechend groß:

\begin{itemize}
    \item \textbf{Kosteneinsparung:} Aktuelle Studien zeigen, dass FM-Kosten
          bis zu 30\,\% durch bessere Planung und Wartungsoptimierung gesenkt werden
          können~\cite{din41011}. Bei einem Marktvolumen von 130 Mrd.\ Euro
          entspräche dies einem Einsparpotenzial von 39 Mrd.\ Euro jährlich allein
          in Deutschland.

    \item \textbf{Energieeffizienz:} Gebäude sind für ca.\ 40\,\% des globalen
          Energieverbrauchs verantwortlich. Optimale Energieplanung durch den
          Subgraph Algorithmus (vgl.\ Abschnitt~8.2.1) könnte den Gebäudeenergiebedarf
          signifikant senken und zum Erreichen der EU-Klimaziele 2030 beitragen.

    \item \textbf{Nachhaltigkeit:} Die polynomielle Lösung des Asset-Scheduling-Problems
          (Abschnitt~\ref{sec:asset}) optimiert Wartungsintervalle so, dass Assets
          ihre volle Lebensdauer ausschöpfen, was Material- und Ressourcenverbrauch
          reduziert.

    \item \textbf{Digitalisierung:} Die ISFMP ist vollständig digital und
          schnittstellen-kompatibel mit BIM (ISO 19650), ERP-Systemen (SAP PM)
          und CAFM-Software (z.\,B.\ ARCHIBUS, Planon). Sie bildet die algorithmische
          Grundlage für KI-gestütztes FM der nächsten Generation.
\end{itemize}

Die Kombination aus theoretischer Fundierung ($P = NP$, Theorem~\ref{thm:pnp}),
praktischer Effizienz ($O(n^3)$, Theorem~\ref{thm:subgraph}) und breiter
Anwendbarkeit (alle FM-Kernprobleme, Proposition~\ref{thm:universal}) macht den
Subgraph Algorithmus zu einem paradigmatischen Werkzeug für die Weiterentwicklung
des Facility Managements als Wissenschaftsdisziplin und Industriepraxis.

\begin{thebibliography}{99}

\bibitem{epp2026subgraph}
    Epp, S. (2026).
    \textit{Der Subgraph Algorithmus}.
    Verfügbar unter: \url{https://github.com/hjstephan86/science}.

\bibitem{garey1979}
    Garey, M.\ R., \& Johnson, D.\ S. (1979).
    \textit{Computers and Intractability: A Guide to the Theory of NP-Completeness}.
    W.\,H.\ Freeman.

\bibitem{cook1971}
    Cook, S.\ A. (1971).
    The complexity of theorem-proving procedures.
    \textit{Proceedings of the 3rd ACM Symposium on Theory of Computing (STOC)},
    151--158.

\bibitem{karp1972}
    Karp, R.\ M. (1972).
    Reducibility among combinatorial problems.
    In R.\ E.\ Miller \& J.\ W.\ Thatcher (Hrsg.),
    \textit{Complexity of Computer Computations}, 85--103.
    Plenum Press.

\bibitem{horrocks2003}
    Horrocks, I., Patel-Schneider, P.\ F., \& van Harmelen, F. (2003).
    From \textsc{SHIQ} and RDF to OWL: The making of a Web Ontology Language.
    \textit{Journal of Web Semantics}, 1(1), 7--26.

\bibitem{din41011}
    Deutsches Institut für Normung. (2019).
    \textit{DIN EN ISO 41011: Facility Management -- Begriffe}.
    Beuth Verlag.

\bibitem{abboud2015}
    Abboud, A., Backurs, A., \& Vassilevska Williams, V. (2015).
    Tight hardness results for LCS and other sequence similarity measures.
    \textit{FOCS 2015}, 59--78.

\bibitem{papadimitriou1994}
    Papadimitriou, C.\ H. (1994).
    \textit{Computational Complexity}.
    Addison-Wesley.

\bibitem{schaefer2001}
    Schaefer, T., \& Umans, C. (2002).
    Completeness in the polynomial-time hierarchy: A compendium.
    \textit{SIGACT News}, 33(3), 32--49.

\bibitem{blum1992}
    Blum, A., \& Rivest, R.\ L. (1992).
    Training a 3-node neural network is NP-complete.
    \textit{Neural Networks}, 5(1), 117--127.

\bibitem{chickering1996}
    Chickering, D.\ M. (1996).
    Learning Bayesian networks is NP-complete.
    In D.\ Fisher \& H.\ Lenz (Hrsg.),
    \textit{Learning from Data}, 121--130.
    Springer.

\bibitem{grover1996}
    Grover, L.\ K. (1996).
    A fast quantum mechanical algorithm for database search.
    \textit{Proceedings of STOC 1996}, 212--219.

\bibitem{kushilevitz1997}
    Kushilevitz, E., \& Nisan, N. (1997).
    \textit{Communication Complexity}.
    Cambridge University Press.

\bibitem{iso19650}
    International Organization for Standardization. (2018).
    \textit{ISO 19650: Organization and digitization of information about buildings
    and civil engineering works, including building information modelling (BIM)}.
    ISO.

\bibitem{nist2024}
    National Institute of Standards and Technology. (2024).
    \textit{Post-Quantum Cryptography Standards: FIPS 203, 204, 205}.
    U.\,S.\ Department of Commerce.

\end{thebibliography}

\end{document}
